

\SGARefTarget{sga3:ix:section:2}\section*{2. Extension of Certain Properties of Diagonalizable Groups to Groups of Multiplicative Type}
\addcontentsline{toc}{section}{2. Extension of Certain Properties of Diagonalizable Groups to Groups of Multiplicative Type}

The extensions in question are essentially trivial consequences of the
results of the preceding Expos\'e, in view of
\SGARefLink{sga3:ix:definition:1:1}{the definitions in 1.1} and the
``local'' nature (for the faithfully flat and quasi-compact topology) of
the results under consideration.

\medskip
\noindent\SGARefTarget{sga3:ix:definition:2:0}\textbf{Definition 2.0.}
\footnote{Editor's note: This definition has been added; it occurs in
\SGARefLink{sga3:ix:proposition:2:3}{Propositions~2.3} and
\SGARefLink{sga3:ix:proposition:2:7}{2.7}.}
\textit{We denote by \(\mathcal M\) the set of faithfully flat and
quasi-compact morphisms.}

\medskip
\noindent\SGARefTarget{sga3:ix:proposition:2:1}\textbf{Proposition 2.1.}
\textit{Let \(G\) be a group of multiplicative type over a prescheme
\(S\). Then:}
\begin{enumerate}
  \item[(a)]\SGARefTarget{sga3:ix:proposition:2:1:item:a:01}
  \textit{\(G\) is faithfully flat over \(S\), and affine over \(S\)
  (a fortiori, quasi-compact over \(S\));}

  \item[(b)]\SGARefTarget{sga3:ix:proposition:2:1:item:b:01}
  \textit{\(G\) is of finite type over \(S\) if and only if \(G\) is of
  finite presentation over \(S\), if and only if, for every \(s\in S\),
  the type of \(G\) at \(s\) is given by a commutative group of finite
  type;}

  \item[(c)]\SGARefTarget{sga3:ix:proposition:2:1:item:c:01}
  \textit{\(G\) is finite over \(S\) if and only if, for every
  \(s\in S\), the type of \(G\) at \(s\) is given by a finite commutative
  group, if and only if (when \(S\) is quasi-compact) \(G\) is of finite
  type over \(S\) and is annihilated by an integer \(n>0\);}

  \item[\((c')\)]\SGARefTarget{sga3:ix:proposition:2:1:item:cprime:01}
  \textit{\(G\) is integral over \(S\) if and only if, for every
  \(s\in S\), the type of \(G\) at \(s\) is given by a torsion
  commutative group;}

  \item[(d)]\SGARefTarget{sga3:ix:proposition:2:1:item:d:01}
  \textit{\(G\) is the unit group if and only if, for every \(s\in S\),
  the type of \(G\) at \(s\) is given by the zero commutative group;}

  \item[(e)]\SGARefTarget{sga3:ix:proposition:2:1:item:e:01}
  \textit{\(G\) is smooth over \(S\) if and only if, for every
  \(s\in S\), the type of \(G\) at \(s\) is given by a commutative group
  of finite type whose torsion subgroup has order prime to the
  characteristic of \(\kappa(s)\).}
\end{enumerate}

These statements follow from \SGARefLink{sga3:viii:proposition:2:1}{VIII~2.1}, taking into account that the
properties in question descend under faithfully flat and quasi-compact
morphisms (cf. SGA~1, VIII, or EGA~IV\textsubscript{2}, \S2).

Using \SGARefLink{sga3:viii:corollary:3:5}{VIII~3.5}, one likewise obtains:

\medskip
\noindent\SGARefTarget{sga3:ix:proposition:2:2}\textbf{Proposition 2.2.}
\textit{Let \(G\) be a group of multiplicative type and of finite type over
\(S\). For every integer \(n\ne0\), the kernel
\({}_nG=\ker(n\cdot\operatorname{id}_G)\) is a group of multiplicative
type, finite over \(S\).}

\medskip
\noindent\SGARefTarget{sga3:ix:proposition:2:3}\textbf{Proposition 2.3.}
\textit{Let \(G\) be a group of multiplicative type over the prescheme
\(S\), acting freely on the right on an \(S\)-prescheme \(X\) that is
affine over \(S\). Then:}
\begin{enumerate}
  \item[(a)]\SGARefTarget{sga3:ix:proposition:2:3:item:a:01}
  \textit{the equivalence relation defined by \(G\) in \(X\) is
  \(\mathcal M\)-effective (\SGARefLink{sga3:iv:subsection:3:4}{IV~3.4}), and \(Y=X/G\) is affine over \(S\);}

  \item[(b)]\SGARefTarget{sga3:ix:proposition:2:3:item:b:01}
  \textit{if, moreover, \(X\) is of finite presentation (respectively, of
  finite type) over \(S\), then the same is true of \(Y\).}
\end{enumerate}

\SGARefLink{sga3:ix:proposition:2:3:item:a:01}{The first assertion} follows from \SGARefLink{sga3:viii:theorem:5:1}{VIII~5.1}, which treats the case in which
\(G\) is diagonalizable, and from \SGARefLink{sga3:iv:proposition:3:5:2}{IV~3.5.2}, which allows one to reduce to
that case. Indeed, faithfully flat and quasi-compact morphisms
\(S'\to S\) are effective descent morphisms for the fibered category of
affine schemes over schemes: if \(Y'\) is affine over \(S'\) and is
equipped with a descent datum relative to \(S'\to S\), then this datum is
effective, that is, \(Y'\) comes from a scheme \(Y\) affine over \(S\)
(cf. SGA~1, \SGARefLink{sga3:viii:proposition:2:1}{VIII~2.1}).

For \SGARefLink{sga3:ix:proposition:2:3:item:b:01}{the second assertion}, one is likewise reduced to the diagonalizable
case of \SGARefLink{sga3:viii:corollary:5:8}{VIII~5.8}, since the finiteness conditions under consideration
descend under faithfully flat and quasi-compact morphisms (SGA~1,
VIII~3.3 and 3.6). Proceeding as in VIII,
Corollaries~5.5--5.7, one deduces from
\SGARefLink{sga3:ix:proposition:2:3}{2.3}:

\medskip
\noindent\SGARefTarget{sga3:ix:corollary:2:4}\textbf{Corollary 2.4.}
\textit{Under the hypotheses of \SGARefLink{sga3:ix:proposition:2:3}{2.3},
the graph morphism}
\SGARefTarget{sga3:ix:corollary:2:4:diagram:01}
\[
  X\times_S G\longrightarrow X\times_S X
\]
\textit{is a closed immersion. For every section \(\sigma\) of \(X\)
over \(S\), the corresponding morphism}
\[
  G\longrightarrow X,\qquad g\longmapsto\sigma\cdot g,
\]
\textit{is a closed immersion.}

In particular:

\medskip
\noindent\SGARefTarget{sga3:ix:corollary:2:5}\textbf{Corollary 2.5.}
\textit{Let \(u:G\to H\) be a monomorphism of \(S\)-preschemes in groups,
with \(G\) of multiplicative type and \(H\) affine over \(S\). Then \(u\)
is a closed immersion, \(H/G=Y\) exists and is affine over \(S\), and
\(H\) is a principal homogeneous bundle over \(Y\) under \(G_Y\).}

\medskip
\noindent\SGARefTarget{sga3:ix:remark:2:6}\textbf{Remark 2.6.}
Let \(u:G\to H\) be a monomorphism of \(S\)-preschemes in groups, where
\(G\) is of multiplicative type and of finite type over \(S\), and \(H\)
is of finite presentation and separated over \(S\). One can then show
that \(u\) is a closed immersion by using \SGARefLink{sga3:viii:proposition:7:12}{VIII~7.12} (see
Remark~\SGARefLink{sga3:viii:remark:7-13:item:b}{VIII~7.13(b)}).

\medskip
\noindent\SGARefTarget{sga3:ix:proposition:2:7}\textbf{Proposition 2.7.}
\textit{Let \(u:G\to H\) be a homomorphism of \(S\)-groups of
multiplicative type, with \(H\) of finite type over \(S\). Then:}
\begin{enumerate}
  \item[(i)]\SGARefTarget{sga3:ix:proposition:2:7:item:i:01}
  \textit{\(G_0=\ker u\) is an \(S\)-group of multiplicative type; the
  equivalence relation defined by \(G_0\) in \(G\) is
  \(\mathcal M\)-effective, and hence (\SGARefLink{sga3:iv:subsection:3:4}{IV~3.4}) \(u\) factors as}
  \SGARefTarget{sga3:ix:proposition:2:7:diagram:01}
  \[
    G\longrightarrow G/G_0=I\longrightarrow H,
  \]
  \textit{where \(I\to H\) is a monomorphism of \(S\)-groups. The group \(I\) is of
  multiplicative type, the equivalence relation in \(H\) defined by \(I\)
  is \(\mathcal M\)-effective, and consequently \(H_0=H/I\) exists and is
  also of multiplicative type;}

  \item[(ii)]\SGARefTarget{sga3:ix:proposition:2:7:item:ii:01}
  \textit{\(H_0\) and \(I\) are of finite type, and so is \(G_0\) when
  \(G\) is.}
\end{enumerate}

\emph{Proof.}
Proceeding as for \SGARefLink{sga3:ix:proposition:2:3}{2.3}, one is
reduced to the case in which \(G\) and \(H\) are diagonalizable; then
\SGARefLink{sga3:ix:proposition:2:7}{2.7} reduces to \SGARefLink{sga3:viii:corollary:3:4}{VIII~3.4}.

We record the following corollaries.

\medskip
\noindent\SGARefTarget{sga3:ix:corollary:2:8}\textbf{Corollary 2.8.}
\begin{enumerate}
  \item[(a)]\SGARefTarget{sga3:ix:corollary:2:8:item:a:01}
  \textit{Let \(S\) be a prescheme. The category of \(S\)-groups of
  multiplicative type and of finite type is an abelian category.}

  \item[(b)]\SGARefTarget{sga3:ix:corollary:2:8:item:b:01}
  \textit{Let \(u:G\to H\) be a homomorphism in this category. It is a
  monomorphism (respectively, an epimorphism, respectively, an
  isomorphism) in this category if and only if \(u\) is a monomorphism of
  preschemes (respectively, \(u\) is faithfully flat, respectively, \(u\)
  is an isomorphism of preschemes).}
\end{enumerate}

It is enough to note that the product of two \(S\)-groups of
multiplicative type is again of multiplicative type (which is immediate),
and is plainly of finite type over \(S\) if both factors are. The rest of
\SGARefLink{sga3:ix:corollary:2:8}{2.8} follows immediately from
\SGARefLink{sga3:ix:proposition:2:7}{2.7}; the details are left to the
reader.

\medskip
\noindent\SGARefTarget{sga3:ix:corollary:2:9}\textbf{Corollary 2.9.}
\textit{Let \(u:G\to H\) be a homomorphism of \(S\)-groups of
multiplicative type and of finite type. Let \(U\) be the set of points
\(s\in S\) for which \(u_s:G_s\to H_s\) is a monomorphism
(respectively, faithfully flat, respectively, an isomorphism). Then
\(U\) is both open and closed, and the induced homomorphism
\(u_U:G_U\to H_U\) is a monomorphism (respectively, faithfully flat,
respectively, an isomorphism).}

Let \(P\) (respectively, \(Q\)) be the kernel (respectively, cokernel) of
\(u\). By \SGARefLink{sga3:ix:proposition:2:7}{2.7}, \(Q\) exists and
\(P\) and \(Q\) are of multiplicative type; moreover, formation of \(P\)
and \(Q\) commutes with every base change \(S'\to S\), in particular with
the formation of fibers. On the other hand, \(u\) is a monomorphism
(respectively, faithfully flat, respectively, an isomorphism) if and only
if \(P=0\) (respectively, \(Q=0\), respectively, \(P=Q=0\)). These
remarks reduce us to checking the following assertion: if \(R\) is an
\(S\)-group of multiplicative type, then the set \(U\) of points
\(s\in S\) for which \(R_s=0\) is open and closed, and \(R_U=0\). This is
contained in \SGARefLink{sga3:ix:remark:1:4:1:item:a}{Remark~1.4.1(a)}.

\medskip
\noindent\SGARefTarget{sga3:ix:corollary:2:10}\textbf{Corollary 2.10.}
\textit{Let \(G\) be an \(S\)-group of multiplicative type and of finite
type, and let \(H\) and \(H'\) be two subgroups of multiplicative type.}
\begin{enumerate}
  \item[(a)]\SGARefTarget{sga3:ix:corollary:2:10:item:a:01}
  \textit{\(H''=H\cap H'\) is a subgroup of \(G\) of multiplicative
  type.}

  \item[(b)]\SGARefTarget{sga3:ix:corollary:2:10:item:b:01}
  \textit{Let \(U\) be the set of \(s\in S\) for which
  \(H_s\subseteq H'_s\) (respectively, \(H_s=H'_s\)). Then \(U\) is open
  and closed, and \(H_U\subseteq H'_U\) (respectively, \(H_U=H'_U\)).}
\end{enumerate}

Here, of course, the intersection sign in \(H\cap H'\) denotes the
intersection in the functorial sense, namely \(H\times_GH'\), which is
plainly an \(S\)-subgroup of \(G\). Likewise, the inclusion sign
(respectively, the equality sign) denotes the order relation
(respectively, equality) between subfunctors of \(H\), not inclusion
(respectively, equality) of the underlying sets.

Applying \SGARefLink{sga3:ix:proposition:2:7}{2.7} first to the inclusion
morphism \(H\to G\), one finds that \(H\) is of finite type; the same
holds for \(H'\). It then follows from
\SGARefLink{sga3:ix:corollary:2:8}{2.8} that \(H\cap H'\) is of
multiplicative type and of finite type. (Note that the canonical functor
from the category considered in \SGARefLink{sga3:ix:corollary:2:8}{2.8}
to the category \((\mathbf{Sch})/S\) commutes with finite inverse limits.)

Formation of \(H\cap H'\) commutes with base extension, in particular
with formation of the fibers. Moreover,
\(H\subseteq H'\) (respectively, \(H=H'\)) is equivalent to
\(H''=H\) (respectively, to \(H''=H\) and \(H''=H'\)). Consider the
canonical homomorphisms \(H''\to H\) and \(H''\to H'\). Then \(U\) is the
set of points \(s\in S\) for which the induced homomorphism
\(H''_s\to H_s\) is an isomorphism (respectively, for which
\(H''_s\to H_s\) and \(H''_s\to H'_s\) are isomorphisms).

By \SGARefLink{sga3:ix:corollary:2:9}{2.9}, it follows that \(U\) is open
and closed, and that \(H''_U\to H_U\) (respectively,
\(H''_U\to H_U\) and \(H''_U\to H'_U\)) are isomorphisms. This gives the
desired conclusion.

\medskip
\noindent\SGARefTarget{sga3:ix:proposition:2:11}\textbf{Proposition 2.11.}
\textit{Let \(S\) be a prescheme, let \(G\) be an \(S\)-group of
multiplicative type and of finite type over \(S\), let \(H\) be a subgroup
of \(G\) of multiplicative type, and put \(K=G/H\), a quotient group of
\(G\) of multiplicative type.}
\begin{enumerate}
  \item[(i)]\SGARefTarget{sga3:ix:proposition:2:11:item:i:01}
  \textit{Suppose \(G\) is trivial, that is, diagonalizable. There is a
  partition of \(S\) into open and closed subsets \(S_i\) such that, for
  every \(i\), \(H_{S_i}\) and \(K_{S_i}\) are diagonalizable. In
  particular, if \(S\) is connected, \(H\) and \(K\) are
  diagonalizable.}

  \item[(ii)]\SGARefTarget{sga3:ix:proposition:2:11:item:ii:01}
  \textit{\SGARefLink{sga3:ix:proposition:2:11:item:i:01}{The same assertion} holds with ``diagonalizable'' replaced by
  ``isotrivial,'' provided that \(S\) is connected, locally connected, or
  quasi-compact.}

  \item[(iii)]\SGARefTarget{sga3:ix:proposition:2:11:item:iii:01}
  \textit{If \(G\) is locally trivial (respectively, locally isotrivial,
  respectively, quasi-isotrivial), then the same is true of \(K\) and
  \(H\).}
\end{enumerate}

\emph{Proof.}
\SGARefLink{sga3:ix:proposition:2:11:item:i:01}{(i)} By hypothesis, \(G=D_S(M)\), where \(M\) is a commutative group of
finite type. For each quotient group \(M_i\) of \(M\), let
\(H_i=D_S(M_i)\) be the corresponding diagonalizable subgroup of \(G\)
(\SGARefLink{sga3:viii:theorem:3:1}{VIII~3.1}). Let \(S_i\) be the set of points \(s\in S\) for which
\(H_s=(H_i)_s\). By \SGARefLink{sga3:ix:corollary:2:10}{2.10}, \(S_i\)
is open and closed, and \(H_{S_i}=(H_i)_{S_i}\); thus \(H_{S_i}\) is
diagonalizable, and hence so is \(K_{S_i}\) (\SGARefLink{sga3:viii:theorem:3:1}{VIII~3.1}). The \(S_i\) are
plainly pairwise disjoint, and we claim that they cover \(S\). Indeed, for
every \(s\), the group \(H_s\), being a subgroup of the diagonalizable
group \(G_s\), is diagonalizable (cf. \SGARefLink{sga3:ix:proposition:8:1}{8.1} below).
Hence \(H_s\) has the form \(D_{\kappa(s)}(M_i)\), by \SGARefLink{sga3:viii:corollary:1:5}{VIII~1.5} and
\SGARefLink{sga3:ix:theorem:3:2:item:b}{3.2(b)}. Restricting to the family of nonempty \(S_i\), assertion \SGARefLink{sga3:ix:proposition:2:11:item:i:01}{(i)}
follows.

\SGARefLink{sga3:ix:proposition:2:11:item:ii:01}{(ii)} By hypothesis, there exists a finite surjective \'etale morphism
\(S'\to S\) such that \(G_{S'}\) is diagonalizable. Every point of \(S'\)
therefore has an open and closed neighborhood \(U'\) such that \(H_{U'}\)
and \(K_{U'}\) are diagonalizable. The image \(U\) of \(U'\) in \(S\) is
open and closed, and \(S'\to S\) still induces a finite surjective
\'etale
morphism \(U'\to U\). Thus every point \(s\in S\) has an open and closed
neighborhood \(U\) such that \(H_U\) and \(K_U\) are isotrivial.
Assertion \SGARefLink{sga3:ix:proposition:2:11:item:ii:01}{(ii)} follows at once.

\SGARefLink{sga3:ix:proposition:2:11:item:iii:01}{(iii)} This follows at once from \SGARefLink{sga3:ix:proposition:2:11:item:i:01}{(i)}, \SGARefLink{sga3:ix:proposition:2:11:item:ii:01}{(ii)}, and the definitions. Moreover,
the quasi-isotrivial case will also follow from the more general fact
``finite type \(\Longrightarrow\) quasi-isotrivial,'' announced in
\SGARefLink{sga3:ix:comments:1:2}{1.2}
(cf. \SGARefLink{sga3:x:corollary:4-5}{X~4.5}).

\SGARefTarget{sga3:ix:section:3}\section*{3. Infinitesimal Properties: Lifting and Conjugation Theorems}
\addcontentsline{toc}{section}{3. Infinitesimal Properties: Lifting and Conjugation Theorems}

The fundamental infinitesimal properties of groups of multiplicative type
follow from \SGARefLink{sga3:ix:theorem:3:1}{the following theorem}.

\medskip
\noindent\SGARefTarget{sga3:ix:theorem:3:1}\textbf{Theorem 3.1.}
\textit{Let \(S\) be an affine scheme, let \(H\) be an \(S\)-group of
multiplicative type, and let \(\mathscr F\) be a quasi-coherent sheaf on
\(S\) on which \(H\) acts (\SGARefLink{sga3:i:subsection:4-7}{I~4.7}). Then}
\[
  \mathrm H^i(H,\mathscr F)=0\qquad\text{for }i>0,
\]
\textit{where \(\mathrm H^i\) denotes the Hochschild cohomology studied
in Expos\'e~I, \S5.}

Indeed, by loc.\ cit., \SGARefLink{sga3:ix:corollary:5:3}{5.3}, Hochschild cohomology is expressed, in terms
of the affine rings \(A\) and \(B\) of \(S\) and \(G\), and the
\(A\)-module \(M\) defining \(\mathscr F\), as the cohomology of a complex
of \(A\)-modules \(C^\bullet(H,M)\), whose formation commutes with every
base change \(A\to A'\). Consequently, for a base change \(A\to A'\) with
\(A'\) flat over \(A\), one has
\[
  \mathrm H^i(H',\mathscr F')
  =\mathrm H^i(H,\mathscr F)\otimes_AA'.
\]
If \(A'\) is moreover faithfully flat over \(A\), it is therefore enough,
in order to prove that \(\mathrm H^i(H,\mathscr F)\) vanishes, to prove
that \(\mathrm H^i(H',\mathscr F')\) vanishes. This observation reduces
us to checking \SGARefLink{sga3:ix:theorem:3:1}{3.1} in the case where
\(G\) is diagonalizable; in that case it was proved in \SGARefLink{sga3:i:theorem:5-3-3}{I~5.3.3}.

Using the results of Expos\'e~III, we shall derive from
\SGARefLink{sga3:ix:theorem:3:1}{3.1} several consequences of a geometric
nature.

\medskip
\noindent\SGARefTarget{sga3:ix:theorem:3:2}\textbf{Theorem 3.2.}
\textit{Let \(S\) be an affine scheme, let \(S_0\) be an affine subscheme
defined by an ideal \(J\), let \(H\) be a group of multiplicative type
over \(S\), and let \(G\) be an arbitrary group prescheme over \(S\).
Let \(u,v:H\to G\) be two group homomorphisms, and let
\(u_0,v_0:H_0\to G_0\) be the homomorphisms obtained from \(u,v\) by the
base change \(S_0\to S\). If \(J^2=0\) and \(u_0=v_0\), then there exists
\(g\in G(S)\) such that}
\[
  v=\operatorname{int}(g)\circ u
  \qquad\textit{and}\qquad
  g_0=e.
\]

This follows from \SGARefLink{sga3:iii:theorem:2-1:item:ii}{III~2.1(ii)} and
\SGARefLink{sga3:ix:theorem:3:1}{3.1} (the vanishing of
\(\mathrm H^1\)).

\medskip
\noindent\SGARefTarget{sga3:ix:corollary:3:3}\textbf{Corollary 3.3.}
\textit{Under the preliminary hypotheses of
\SGARefLink{sga3:ix:theorem:3:2}{3.2}, suppose in addition that \(G\) is
smooth over \(S\), but assume only that \(J\) is nilpotent (not
necessarily square-zero). If there exists \(g_0\in G_0(S_0)\) such that
\(v_0=\operatorname{int}(g_0)\circ u_0\), then \(g_0\) lifts to an element
\(g\in G(S)\) such that \(v=\operatorname{int}(g)\circ u\).}

By induction on an integer \(n\) such that \(J^n=0\), one is reduced to
the case in which \(J\) is square-zero. Since \(G\) is smooth over \(S\),
one may lift \(g_0\) to an element \(g'\) of \(G(S)\). Replacing \(v\) by
\(v'=\operatorname{int}(g')\circ u\), one is then reduced to the
situation of \SGARefLink{sga3:ix:theorem:3:2}{3.2}.

\medskip
\noindent\SGARefTarget{sga3:ix:corollary:3:4}\textbf{Corollary 3.4.}
\textit{Under the preliminary hypotheses of
\SGARefLink{sga3:ix:theorem:3:2}{3.2}, with \(J\) nilpotent, suppose in
addition that \(u\) is central (for example, that \(G\) is commutative).
Then \(u_0=v_0\) implies \(u=v\).}

Indeed, the reduction to the case in which \(J\) is square-zero is again
immediate, and then it is enough to apply
\SGARefLink{sga3:ix:theorem:3:2}{3.2}. In particular:

\medskip
\noindent\SGARefTarget{sga3:ix:corollary:3:5}\textbf{Corollary 3.5.}
\textit{Let \(S,S_0,H,G\) be as in
\SGARefLink{sga3:ix:theorem:3:2}{3.2}, with \(J\) nilpotent. Let
\(u:H\to G\) be a homomorphism of \(S\)-groups such that
\(u_0:H_0\to G_0\) is the unit homomorphism. Then \(u\) is the unit
homomorphism.}

\medskip
\noindent\SGARefTarget{sga3:ix:theorem:3:6}\textbf{Theorem 3.6.}
\textit{Let \(S\) be an affine scheme, let \(S_0\) be a subscheme defined
by a nilpotent ideal \(J\), let \(H\) be a group of multiplicative type
over \(S\), let \(G\) be a group prescheme smooth over \(S\), and let
\(u_0:H_0\to G_0\) be a homomorphism of the \(S_0\)-groups obtained by
base change \(S_0\to S\). Then there exists a homomorphism
\(u:H\to G\) of \(S\)-groups that lifts \(u_0\). Moreover, by
\SGARefLink{sga3:ix:corollary:3:3}{3.3}, any two such liftings \(u,u'\)
are conjugate by an element of \(G(S)\) whose reduction in \(G(S_0)\) is
the identity element.}

This follows from \SGARefLink{sga3:iii:theorem:2-1}{III~2.1} and \SGARefLink{sga3:ix:proposition:2:3}{2.3} and from
\SGARefLink{sga3:ix:theorem:3:1}{3.1}: the vanishing of
\(\mathrm H^2\) gives the existence of the lifting of \(u_0\), and the
vanishing of \(\mathrm H^1\) gives uniqueness up to conjugation.

One may prove in the same way the following variants of
\SGARefLink{sga3:ix:theorem:3:2}{3.2}--%
\SGARefLink{sga3:ix:theorem:3:6}{3.6}. In fact, these variants imply the
preceding statements when applied to the graph subgroups of the
homomorphisms considered there.

\medskip
\noindent\SGARefTarget{sga3:ix:theorem:3:2:bis}\textbf{Theorem 3.2 bis.}
\textit{Let \(S\) be an affine scheme, let \(S_0\) be a subscheme defined
by an ideal \(J\), let \(G\) be an \(S\)-group prescheme, and let \(H,H'\)
be subgroup preschemes of multiplicative type. Let \(G_0,H_0,H'_0\) be
the groups over \(S_0\) obtained by base change \(S_0\to S\). If
\(J^2=0\) and \(H_0=H'_0\), then there exists \(g\in G(S)\) such that}
\[
  H'=\operatorname{int}(g)(H)
  \qquad\textit{and}\qquad
  g_0=e.
\]

\medskip
\noindent\SGARefTarget{sga3:ix:corollary:3:3:bis}\textbf{Corollary 3.3 bis.}
\textit{Under the preliminary hypotheses of
\SGARefLink{sga3:ix:theorem:3:2:bis}{3.2 bis}, suppose in addition that
\(G\) is smooth over \(S\), and that \(J\) is nilpotent (not necessarily
square-zero). If \(g_0\in G_0(S_0)\) satisfies
\(H'_0=\operatorname{int}(g_0)(H_0)\), then \(g_0\) lifts to an element
\(g\in G(S)\) such that \(H'=\operatorname{int}(g)(H)\).}

\medskip
\noindent\SGARefTarget{sga3:ix:corollary:3:4:bis}\textbf{Corollary 3.4 bis.}
\textit{Under the preliminary hypotheses of
\SGARefLink{sga3:ix:theorem:3:2:bis}{3.2 bis}, suppose that \(J\) is
nilpotent and that \(H\) is central (for example, that \(G\) is
commutative). Then \(H_0=H'_0\) implies \(H=H'\).}

\medskip
\noindent\SGARefTarget{sga3:ix:theorem:3:6:bis}\textbf{Theorem 3.6 bis.}
\textit{Let \(S\) be an affine scheme, let \(S_0\) be a subscheme defined
by a nilpotent ideal \(J\), let \(G\) be an \(S\)-group prescheme smooth
over \(S\), and let \(H_0\) be a subgroup prescheme of
\(G_0=G\times_SS_0\), of multiplicative type. Then:}
\begin{enumerate}
  \item[(a)]\SGARefTarget{sga3:ix:theorem:3:6:bis:item:a:01}
  \textit{there exists a subgroup prescheme \(H\) of \(G\), flat over
  \(S\), such that \(H\times_SS_0=H_0\);}

  \item[\SGAInlineRefTarget{sga3:ix:theorem:3:2:item:b}(b)]\SGARefTarget{sga3:ix:theorem:3:6:bis:item:b:01}
  \textit{such an \(H\) is necessarily of multiplicative type;}

  \item[(c)]\SGARefTarget{sga3:ix:theorem:3:6:bis:item:c:01}
  \textit{finally, by
  \SGARefLink{sga3:ix:theorem:3:2:bis}{3.2 bis}, any two such liftings
  \(H,H'\) of \(H_0\) are conjugate by an element \(g\in G(S)\) whose
  reduction in \(G(S_0)\) is the identity element.}
\end{enumerate}

Let us indicate how one can prove
\SGARefLink{sga3:ix:theorem:3:2:bis}{3.2 bis} (of which
\SGARefLink{sga3:ix:corollary:3:3:bis}{3.3 bis} and
\SGARefLink{sga3:ix:corollary:3:4:bis}{3.4 bis} are immediate
consequences) and assertion
\SGARefLink{sga3:ix:theorem:3:6:bis:item:a:01}{(a)} of
\SGARefLink{sga3:ix:theorem:3:6:bis}{3.6 bis}. The latter follows from
\SGARefLink{sga3:ix:theorem:3:1}{3.1} (the vanishing of
\(\mathrm H^2\)) and \SGARefLink{sga3:iii:corollary:4-37:item:ii}{III~4.37(ii)}. Similarly,
\SGARefLink{sga3:ix:theorem:3:2:bis}{3.2 bis} follows from
\SGARefLink{sga3:ix:theorem:3:1}{3.1} (the vanishing of
\(\mathrm H^1\)) and \SGARefLink{sga3:iii:corollary:4-37:item:i}{III~4.37(i)}, at least when \(G\) and \(H\) are
smooth over \(S\).

Using the results of the following Expos\'e, however, one can very simply
deduce \SGARefLink{sga3:ix:theorem:3:2:bis}{3.2 bis} and
\SGARefLink{sga3:ix:theorem:3:6:bis}{3.6 bis}, in the general form stated
here, from \SGARefLink{sga3:ix:theorem:3:2}{3.2} and
\SGARefLink{sga3:ix:theorem:3:6}{3.6}, respectively. For
\SGARefLink{sga3:ix:theorem:3:2:bis}{3.2 bis}, it is enough to note that,
by \SGARefLink{sga3:x:proposition:2:1}{X~2.1}, \(H\) and \(H'\) are isomorphic, which reduces us to
\SGARefLink{sga3:ix:theorem:3:2}{3.2}.

For \SGARefLink{sga3:ix:theorem:3:6:bis}{3.6 bis}, note that \(H_0\) is
necessarily of finite type over \(S_0\):\footnote{Editor's note: The
original said, ``by \SGARefLink{sga3:ix:proposition:2:1:item:b:01}{2.1(b)}, since its fibers are.'' The editors did not
understand this argument and have replaced it by the argument that
follows.} indeed, as a subgroup prescheme of \(G_0\), which is smooth
over \(S_0\), the group \(H_0\) is locally of finite type over \(S_0\).
But \(H_0\) is affine over \(S_0\) by
\SGARefLink{sga3:ix:proposition:2:1:item:a:01}{2.1(a)}, and is therefore
of finite type over \(S_0\). By \SGARefLink{sga3:x:corollary:4-5}{X~4.5}, \(H_0\) is consequently
quasi-isotrivial, and hence comes, by \SGARefLink{sga3:x:proposition:2:1}{X~2.1}, from a group \(H\) of
multiplicative type over \(S\).\footnote{Editor's note: The sentence that
follows has been expanded.} By
\SGARefLink{sga3:ix:theorem:3:6}{3.6}, there exists a homomorphism of
\(S\)-groups \(u:H\to G\) lifting the immersion \(H_0\hookrightarrow
G_0\). Since \(H\) and \(H_0\) (respectively, \(G\) and \(G_0\)) have the
same underlying topological space, and since, for every \(h\in H\), the
morphism

\[
  \mathscr O_{G,u(h)}\longrightarrow\mathscr O_{H,h}
\]
is surjective (because it is so after reduction modulo the nilpotent ideal
\(J\)), the morphism \(u\) is also an immersion (cf. EGA~I, 4.2.2).

Finally, every lifting \(H\) of \(H_0\) to a flat subgroup of \(G\) is
necessarily of multiplicative type by \SGARefLink{sga3:x:corollary:2:3}{X~2.3}, which
also proves assertion
\SGARefLink{sga3:ix:theorem:3:6:bis:item:b:01}{(b)} of
\SGARefLink{sga3:ix:theorem:3:6:bis}{3.6 bis}. The reader may verify that
the results \SGARefLink{sga3:ix:theorem:3:2:bis}{3.2 bis}--%
\SGARefLink{sga3:ix:theorem:3:6:bis}{3.6 bis} are not used in
Expos\'e~X,
so there is no vicious circle.

\medskip
\noindent\SGARefTarget{sga3:ix:proposition:3:8}\textbf{Proposition 3.8.}
\textit{Let \(S\) be a prescheme, let \(S_0\) be a closed subscheme
defined by a locally nilpotent ideal \(J\), let \(G\) be an \(S\)-group
of multiplicative type, and let \(X\) be an \(S\)-prescheme locally of
finite presentation. Suppose that \(G\) acts on \(X\) in such a way that
\(G_0=G\times_SS_0\) acts trivially on \(X_0=X\times_SS_0\). Then \(G\)
acts trivially on \(X\).}

The proof is the proof of \SGARefLink{sga3:iii:corollary:2-4:item:b}{III~2.4(b)}. One may also reduce to loc.\ cit. by
observing that \SGARefLink{sga3:ix:proposition:3:8}{3.8} immediately
reduces to the case in which \(G\) is diagonalizable, which is the case
treated there.

\medskip
\noindent\SGARefTarget{sga3:ix:corollary:3:9}\textbf{Corollary 3.9.}
\textit{Let \(S,S_0\) be as above, and let \(u:G\to H\) be a homomorphism
of \(S\)-groups, with \(G\) of multiplicative type and \(H\) locally of
finite presentation over \(S\). Suppose that the homomorphism
\(u_0:G_0\to H_0\) obtained by base change \(S_0\to S\) is central. Then
\(u\) is central.}

Indeed, it is enough to apply
\SGARefLink{sga3:ix:proposition:3:8}{3.8} with \(X=H\), letting \(G\) act
on \(H\) by
\[
  (g,h)\longmapsto\operatorname{int}(u(g))\cdot h
  =u(g)h\,u(g)^{-1}.
\]

